\subsection{卢津空间与波莱尔集}

定理 3. 设 X 是卢津空间。则 X 的子空间是卢津空间，当且仅当它是波莱尔集。

此定理是下述两个引理的推论：

引理 6. 在卢津空间 X 中，每个波莱尔集都是 X 的卢津子空间。

设 \(\mathfrak{X}\) 是由 X 的诸子集 A 中 A 与 \(\complement A\) 都是 X 的卢津子空间者所成的集。由于 X 中每个闭集与每个开集都是 X 的卢津子空间（第 4 号），\(\mathfrak{X}\) 包含 X 的所有闭子集，因此只要证明 \(\mathfrak{X}\) 是 X 上的一个 \(\sigma\) 代数，引理即得证。为此只需证明：若 \((A_n)\) 是 \(\mathfrak{X}\) 中集合的一个序列，则 \(\bigcap_{n} A_n\) 与 \(\bigcup_{n} A_n\) 都是 X 的卢津子空间。我们在第 4 号中已看到，卢津子空间的每个可数交是卢津子空间。另一方面，若 \(B_n\) 是 \(A_n\) 与 \(\bigcap_{i < n} \complement A_i\) 的交，则由假设与前面的注，\(B_n\) 是卢津子空间；且由于 \(\bigcup_{n} A_n = \bigcup_{n} B_n\) ，由第 4 号的引理 2 ，子空间 \(\bigcup_{n} A_n\) 是卢津子空间。

引理 7. 可度量化空间 X 的每个卢津子空间 A 都是 X 中的波莱尔集。

由第 5 号的命题 13 ，存在一个筛 C 与一个连续双射 \(f\) （ \(\mathbf{L}(\mathbf{C})\) 到 A 上的）。用第 6 号引理 5 的记号，对每个整数 \(n\) 与每个 \(c \in \mathbf{C}_n\) ，令 \(g_n(c)\) 表示 X 的子空间 \(f(q_n(c))\) ；它是卢津子空间，因而是 X 的苏斯林子空间。当 \(c\) 跑遍 \(\mathbf{C}_n\) 时，诸集 \(g_n(c)\) 两两不相交，因为 \(f\) 是双射；故由第 6 号的定理 2 ，存在 X 中波莱尔集的一个族 \(c \to g_n'(c) (c \in \mathbf{C}_n)\) ，两两不相交，且使 \(g_n'(c) \supset g_n(c)\) 对一切 \(c \in \mathbf{C}_n\) 成立。必要时以 \(g_n'(c)\) 与闭包 \(\overline{g_n(c)}\) （即 \(g_n(c)\) 在 X 中的闭包）的交代替之，可设 \(g_n'(c) \subset \overline{g_n(c)}\) 。令 \(c_{n-1}, c_{n-2}, \ldots, c_0\) 表示 \(c\) 在 \(\mathbf{C}_{n-1}, \mathbf{C}_{n-2}, \ldots, \mathbf{C}_0\) 中的像，所用满射为

\[
p _ {n - 1, n} = p _ {n - 1}, p _ {n - 2, n} = p _ {n - 2} \circ p _ {n - 1}, \dots , p _ {0 n} = p _ {0} \circ p _ {1} \circ \dots \circ p _ {n - 1}
\]

（分别为）；并令 \(h_n(c)\) 表示下列诸集

\[
g _ {n} ^ {\prime} (c), g _ {n - 1} ^ {\prime} \left(c _ {n - 1}\right), \dots , g _ {o} ^ {\prime} \left(c _ {0}\right).
\]

的交。由于 \(q_{i}(c_{i}) \supset q_{n}(c)\) 对 \(0 \leqslant i \leqslant n - 1\) 成立，\(h_{n}(c)\) 包含 \(g_{n}(c)\) ；同样显然 \(h_{n}(c)\) 是波莱尔集且含于 \(g_{n}(c)\) ，且当 \(c\) 跑遍 \(\mathbf{C}_{n}\) 时诸 \(h_{n}(c)\) 两两不相交；最后，由构造，对每个 \(c' \in \mathbf{C}_{n+1}\) 有 \(h_{n+1}(c') \subset h_n(p_n(c'))\) 。令 \(\mathbf{B}_n\) 为诸集 \(h_n(c)\) （当 \(c\) 跑遍 \(\mathbf{C}_n\) ）的并；\(\mathbf{B}_n\) 是波莱尔集，且 \(\mathbf{B}_{n+1} \subset \mathbf{B}_n\) ；又 \(\mathbf{B}_n\) 包含诸集 \(g_n(c) (c \in \mathbf{C}_n)\) 的并，即 A 。令 B 为递减集序列 \(\mathbf{B}_n\) 的交；B 是波莱尔集且包含 A 。我们将证明 B = A ，这样就完成了证明。

设 \(x\) 是 \(\mathbf{B}\) 的一点。则对每个整数 \(n\) ，存在唯一的 \(c \in \mathbf{C}_n\) 使 \(x \in h_n(c)\) ；我们把此 \(c\) 记作 \(c_n(x)\) 。序列 \((c_n(x))_{n \geqslant 0}\) 属于 \(\mathbf{L}(\mathbf{C})\) 。递减序列 \((g_n(c_n(x)))\) 依定义收敛于一点 \(a \in \mathbf{A}\) ；这些集的闭包序列也收敛于 \(a\) （在 \(\mathbf{X}\) 中），故序列 \((h_n(c_n(x)))\) 更加如此。而 \(x\) 属于所有集 \(h_n(c_n(x))\) ，因此 \(x = a \in \mathbf{A}\) 。引理 7 由此得证，定理 3 亦随之得证。

推论. 若 \(f\) 是卢津空间（特别地，波兰空间）\(\mathbf{X}\) 到可度量化空间 \(\mathbf{Y}\) 中的连续单射映射，则 \(f(\mathbf{X})\) 是 \(\mathbf{Y}\) 中的波莱尔集。

